\section{充分必要条件和直线与圆结合}

\begin{tips}
    填空题，是否充分，是否必要？
\end{tips}

\begin{example}
    设 $a$ 为实数，直线 $l_1: a x+y=1$ ，直线 $l_2: x+a y=2 a$ ，则＂$a=1$＂是＂$l_1, l_2$ 平行＂的（\quad）条件
\end{example}

\v{8em}

\begin{example}
    "$m=-4$"是"直线 $l_1:(m-2) x-3 y-1=0$ 与直线 $l_2: m x+(m+2) y+1=0$ 相互平行"的 （\qquad）
\end{example}

\begin{example}
    ＂ $3<r<6$＂是＂圆 $C:(x+1)^2+(y-2)^2=r^2(r>0)$ 上

    恰有 2 个点到直线 $l: 3 x+4 y+15=0$ 的距离为 1 ＂的（\quad ）条件
\end{example}

\v{8em}

\begin{example}
    已知直线 $l_1: a x+2 y=0$ 与直线 $l_2: x+(a+1) y+4=0$ ，则＂$a=1$＂是＂$l_1 / / l_2$＂ （\qquad）
\end{example}

\v{8em}

\begin{example}
    已知直线 $l: x+2 y+t=0$ ，曲线 $C: y=\sqrt{4-x^2}$ ，则＂$l$ 与 $C$ 相切＂是＂$t=2 \sqrt{5}$＂的 （\qquad）
\end{example}

\v{8em}

\begin{example}
    ＂关于 $x, y$ 的方程：$x^2+y^2+m x+4 y+8=0$ 表示圆＂是＂$m>4$＂的（ \quad）条件
\end{example}

\v{8em}

\begin{example}
    已知直线 $a x+\sqrt{3} y+2=0$ 与圆 $x^2+y^2=4$ 交于 $A, B$ 两点，则＂$|A B|=2 \sqrt{3}$＂是＂$a=1$＂的 （\qquad）
\end{example}

\v{8em}

\begin{example}
    ＂$a>0$＂是＂直线 $x-a y+2 a-1=0(a \in \mathbf{R})$ 与圆 $x^2+y^2=1$ 相交＂的 （\qquad）
\end{example}

\v{8em}

\begin{example}
    ＂$a^2+b^2<R^2$＂是＂圆 $(x-a)^2+(y-b)^2=R^2$ 与坐标轴有四个交点＂的 （\qquad）
\end{example}

\v{8em}

\begin{example}
    已知圆 $C_1:(x-2 m)^2+(y-2 m)^2=9(m-2)$ 与圆 $C_2: x^2+y^2-8 x-8 y+34-m=0$

    则＂$m=4$＂是＂圆 $C_1$ 与圆 $C_2$ 外切＂的 （\qquad）
\end{example}

\v{8em}

\subsection{练习参考答案}

\begin{solution}
    充分不必要条件

    \textbf{解答}
    若 $a=1$，则直线 $l_1: x+y=1$，直线 $l_2: x+y=2$。此时两直线平行，所以充分性成立。
    若 $l_1 \parallel l_2$，则有 $a \cdot 1 = 1 \cdot 1$，解得 $a^2=1$，即 $a = \pm 1$。  当 $a=1$ 或 $a=-1$ 时，两直线均平行。  因此，$l_1 \parallel l_2$ 不能唯一推出 $a=1$，必要性不成立。
    故为充分不必要条件。
\end{solution}

\begin{solution}
    充要条件

    \textbf{解答}
    两直线平行的条件是 $(m-2)(m+2) = -3m$，且截距不相等。
    化简得 $m^2+3m-4=0$，解得 $m=-4$ 或 $m=1$。
    当 $m=-4$ 时，两直线分别为 $-6x-3y-1=0$ 和 $-4x-2y+1=0$，斜率均为 $-2$，截距不同，互相平行。
    当 $m=1$ 时，两直线方程相同，为重合关系，不符合平行。
    因此，$m=-4$ 是两直线平行的充要条件。
\end{solution}

\begin{solution}
    必要不充分条件

    \textbf{解答}
    圆上恰有2个点到直线 $l:3x+4y+15=0$ 的距离为1，等价于该圆与到直线 $l$ 距离为1的两条平行线 $3x+4y+10=0$ 和 $3x+4y+20=0$ 恰有两个交点。
    圆心 $C(-1, 2)$ 到这两条直线的距离分别为 $d_1 = \frac{|-3+8+10|}{5}=3$ 和 $d_2 = \frac{|-3+8+20|}{5}=5$。
    为了恰好有两个交点，圆必须与一条直线相交，与另一条直线相离，即半径 $r$ 需满足 $d_1 < r < d_2$，所以 $3 < r < 5$。
    命题 $P: 3 < r < 6$，命题 $Q: 3 < r < 5$。因为 $Q \implies P$ 但 $P \not\implies Q$，所以 $P$ 是 $Q$ 的必要不充分条件。
\end{solution}

\begin{solution}
    充分不必要条件

    \textbf{解答}
    若 $a=1$，则直线 $l_1: x+2y=0$，$l_2: x+2y+4=0$，两直线平行。充分性成立。
    若 $l_1 \parallel l_2$，则 $a(a+1) = 1 \cdot 2$，解得 $a=1$ 或 $a=-2$。
    当 $a=-2$ 时，两直线也平行。  因此，$l_1 \parallel l_2$ 不能唯一推出 $a=1$，必要性不成立。
    故为充分不必要条件。
\end{solution}

\begin{solution}
    既不充分也不必要条件

    \textbf{解答}
    曲线 $C: y=\sqrt{4-x^2}$ 是圆心为 $(0,0)$、半径为 $r=2$ 的上半圆。
    若 $t=2\sqrt{5}$，圆心到直线 $x+2y+2\sqrt{5}=0$ 的距离 $d = \frac{|2\sqrt{5}|}{\sqrt{1^2+2^2}}=2=r$。  但切点 $y<0$，是与下半圆相切，故充分性不成立。
    若直线 $l$ 与上半圆 $C$ 相切，则圆心到直线距离为2，且切点 $y \ge 0$。
    由 $\frac{|t|}{\sqrt{5}}=2$ 得 $t = \pm 2\sqrt{5}$。  经检验，只有当 $t=-2\sqrt{5}$ 时切点位于上半圆。  因此，相切推不出 $t=2\sqrt{5}$，必要性不成立。
    故为既不充分也不必要条件。
\end{solution}

\begin{solution}
    必要不充分条件

    \textbf{解答}
    将方程 $x^2+y^2+mx+4y+8=0$ 配方得 $(x+\frac{m}{2})^2 + (y+2)^2 = \frac{m^2}{4}-4$。
    该方程表示圆的条件是半径的平方大于零，即 $\frac{m^2}{4}-4 > 0$，解得 $m>4$ 或 $m<-4$。
    设 A: "方程表示圆" ($m>4$ 或 $m<-4$)，B: "$m>4$"。
    B $\implies$ A，但 A $\not\implies$ B (例如 $m=-5$)。因此，A 是 B 的必要不充分条件。
\end{solution}

\begin{solution}
    必要不充分条件

    \textbf{解答}
    圆 $x^2+y^2=4$ 的半径 $r=2$。由弦长公式 $|AB|=2\sqrt{r^2-d^2}$，得 $2\sqrt{3}=2\sqrt{4-d^2}$，解得弦心距 $d=1$。
    圆心 $(0,0)$ 到直线 $ax+\sqrt{3}y+2=0$ 的距离 $d = \frac{|2|}{\sqrt{a^2+3}}$。
    令 $d=1$，即 $\frac{2}{\sqrt{a^2+3}}=1$，解得 $a^2=1$，即 $a=\pm 1$。
    设 A: "$|AB|=2\sqrt{3}$" ($a=\pm 1$)，B: "$a=1$"。
    B $\implies$ A，但 A $\not\implies$ B。因此，A 是 B 的必要不充分条件。
\end{solution}

\begin{solution}
    必要不充分条件

    \textbf{解答}
    直线与圆相交的条件是圆心到直线的距离小于半径。圆心为 $(0,0)$，半径为 $1$。
    距离 $d = \frac{|2a-1|}{\sqrt{1+a^2}} < 1$。
    两边平方得 $(2a-1)^2 < 1+a^2$，化简得 $3a^2-4a<0$，解得 $0 < a < \frac{4}{3}$。
    设 A: "$a>0$"，B: "直线与圆相交" ($0<a<\frac{4}{3}$)。
    B $\implies$ A，但 A $\not\implies$ B (例如 $a=2$) 。因此，A 是 B 的必要不充分条件。
\end{solution}

\begin{solution}
    充分不必要条件

    \textbf{解答}
    条件 “$a^2+b^2 < R^2$” 表示圆心 $(a,b)$ 到原点的距离小于半径 $R$，即原点 $(0,0)$ 在圆内部。  若原点在圆内，则圆必与 x 轴和 y 轴均有两个交点，共四个交点。所以充分性成立。
    反之，若圆与坐标轴有四个交点，则圆心到 x 轴和 y 轴的距离都必须小于半径，即 $|b|<R$ 且 $|a|<R$。  但这不能推出 $a^2+b^2 < R^2$ (例如 $a=0.8R, b=0.8R$) 。所以必要性不成立。
    故为充分不必要条件。
\end{solution}

\begin{solution}
    充要条件

    \textbf{解答}
    圆 $C_1$ 的圆心为 $(2m, 2m)$，半径 $r_1=3\sqrt{m-2}$。
    圆 $C_2$ 的标准方程为 $(x-4)^2+(y-4)^2=m-2$，圆心为 $(4,4)$，半径 $r_2=\sqrt{m-2}$。
    两圆外切的条件为圆心距等于半径之和，即 $|C_1C_2| = r_1+r_2$。
    $|C_1C_2| = \sqrt{(2m-4)^2+(2m-4)^2} = 2\sqrt{2}|m-2|$。
    $r_1+r_2=4\sqrt{m-2}$。
    由于 $m>2$，则 $2\sqrt{2}(m-2) = 4\sqrt{m-2}$。  解得 $m=4$ ( $m=2$ 舍去)。
    因此，“$m=4$” 是两圆外切的充要条件。
\end{solution}

\begin{solution}
    必要不充分条件

    \textbf{解答}
    方程表示焦点在 y 轴上的双曲线，需要满足分母一个为正，一个为负，且 $y^2$ 的分母为正。
    即 $7-m>0$ 且 $m-3<0$。
    解得 $m<3$。
    设 A: "$m<3$ 或 $m>7$"，B: "方程为焦点在 y 轴的双曲线" ($m<3$)。
    B $\implies$ A，但 A $\not\implies$ B (例如 $m=8$) 。因此，A 是 B 的必要不充分条件。
\end{solution}